\documentclass[11pt,a4paper]{article}
\usepackage[margin=1in]{geometry}
\usepackage{amsmath,amssymb,amsthm}
\usepackage{algorithm}
\usepackage{algpseudocode}
\usepackage{booktabs}
\usepackage{hyperref}

\theoremstyle{definition}
\newtheorem{definition}{Definition}[section]
\newtheorem{theorem}{Theorem}[section]
\newtheorem{lemma}[theorem]{Lemma}
\newtheorem{remark}{Remark}[section]

\title{\textbf{Theoretical Foundations, Multi-Dimensional Pareto Optimization, and Resource Balance in Compound Model Scaling}}
\author{Team Scaibu \\ \small Email: \texttt{jaydeep.9664920749@gmail.com} \\ \small Architecture and Core Engine Team}
\date{October 2026}

\begin{document}
\maketitle

\begin{abstract}
This document provides the formal mathematical specification, Pareto frontier balance, and dimensional resource allocation analysis of Compound Model Scaling (EfficientNet). Conventional scaling methods arbitrarily expand a single dimension: depth (layer count $d$), width (channel channels $w$), or image resolution ($r$). Compound scaling demonstrates that balancing all three dimensions under a unified compound coefficient $\phi$ achieves optimal accuracy per FLOP by maintaining balanced receptive field and feature capacity. We formulate the constrained optimization problem, derive exponential FLOP scaling relations, detail hardware channel alignment rules, trace a worked numerical example, and document baseline grid search heuristics.
\end{abstract}

\section{Notation and Mathematical Preliminaries}

Let $\mathcal{N}(\theta) = \bigotimes_{i=1}^S \mathcal{F}_i^{d_i}(X_{\langle h_i, w_i, c_i \rangle})$ denote a baseline convolutional network partitioned into $S$ sequential stages.

\begin{itemize}
    \item $S \in \mathbb{N}_{\ge 1}$: Number of distinct architectural stages.
    \item $d = (d_1, \dots, d_S) \in \mathbb{N}^S$: Baseline layer depth repeats per stage.
    \item $c = (c_1, \dots, c_S) \in \mathbb{N}^S$: Baseline channel capacities per stage.
    \item $r_0 \in \mathbb{N}_{\ge 1}$: Baseline input image spatial resolution ($r_0 = 224$).
    \item $\phi \ge 0$: Global user-specified compound scaling coefficient.
    \item $\alpha \ge 1.0$: Depth scaling base coefficient.
    \item $\beta \ge 1.0$: Width scaling base coefficient.
    \item $\gamma \ge 1.0$: Resolution scaling base coefficient.
    \item $\delta \in \mathbb{N}_{\ge 1}$: Hardware memory alignment divisor ($\delta = 8$ for GPU tensor cores).
\end{itemize}

\begin{table}[h]
\centering
\caption{Summary of Mathematical Symbols for Compound Scaling}
\begin{tabular}{lll}
\toprule
\textbf{Symbol} & \textbf{Type / Domain} & \textbf{Description} \\
\midrule
$\phi$ & $\mathbb{R}_{\ge 0}$ & Compound scale power parameter \\
$\alpha, \beta, \gamma$ & $\mathbb{R}_{\ge 1.0}$ & Coordinate base scale constants \\
$\hat{d}_i$ & $\mathbb{N}_{\ge 1}$ & Scaled depth repeat count for stage $i$ \\
$\hat{c}_i$ & $\mathbb{N}_{\ge 1}$ & Scaled channel count for stage $i$ (multiple of $\delta$) \\
$\hat{r}$ & $\mathbb{N}_{\ge 1}$ & Scaled input image spatial resolution \\
$\text{Ratio}_{\text{FLOPs}}$ & $\mathbb{R}_{\ge 1.0}$ & Theoretical compute scaling multiplier ($\approx 2^\phi$) \\
\bottomrule
\end{tabular}
\end{table}

\section{Problem Definition and Core Concepts}

\subsection{Intuitive Meaning}
Scaling only network depth captures richer and more complex features, but suffers from diminishing returns due to vanishing gradients. Scaling only width captures fine-grained features, but wide shallow networks struggle with complex abstractions. Scaling only image resolution captures finer detail, but resolution gains saturate quickly. Compound scaling scales depth, width, and resolution simultaneously in balanced proportions, ensuring that larger images have deeper networks with wider channels to process larger receptive fields.

\subsection{Formal Constrained Optimization Problem}
The compound scaling formulation optimizes accuracy subject to resource constraints:
\begin{align}
\text{depth: } & d = \alpha^\phi \\
\text{width: } & w = \beta^\phi \\
\text{resolution: } & r = \gamma^\phi \\
\text{subject to: } & \alpha \cdot \beta^2 \cdot \gamma^2 \approx 2 \quad \text{and } \alpha \ge 1, \, \beta \ge 1, \, \gamma \ge 1
\end{align}

\section{Algorithm Specification}

\begin{algorithm}[H]
\caption{Compound Architecture Scaling Execution}
\label{alg:compound_scaling}
\begin{algorithmic}[1]
\Require Base depths $d \in \mathbb{N}^S$, base channels $c \in \mathbb{N}^S$, base resolution $r_0$, compound factor $\phi \ge 0$, base constants $(\alpha, \beta, \gamma)$, divisor $\delta$.
\Ensure Scaled depths $\hat{d}$, scaled channels $\hat{c}$, scaled resolution $\hat{r}$, FLOP ratio $\rho$.
\State $d_{\text{mult}} \gets \alpha^\phi$, $w_{\text{mult}} \gets \beta^\phi$, $r_{\text{mult}} \gets \gamma^\phi$
\State Allocate $\hat{d} \in \mathbb{N}^S$, $\hat{c} \in \mathbb{N}^S$
\For{$i \gets 1 \textbf{ to } S$}
    \State $\hat{d}[i] \gets \max(1, \lceil d[i] \cdot d_{\text{mult}} \rceil)$
    \State $c_{\text{raw}} \gets c[i] \cdot w_{\text{mult}}$
    \State $c_{\text{round}} \gets \max(\delta, \lfloor (c_{\text{raw}} + \delta / 2) / \delta \rfloor \cdot \delta)$
    \If{$c_{\text{round}} < 0.9 \cdot c_{\text{raw}}$}
        \State $c_{\text{round}} \gets c_{\text{round}} + \delta$
    \EndIf
    \State $\hat{c}[i] \gets c_{\text{round}}$
\EndFor
\State $\hat{r} \gets \text{round}(r_0 \cdot r_{\text{mult}})$
\State $\rho \gets d_{\text{mult}} \cdot (w_{\text{mult}})^2 \cdot (r_{\text{mult}})^2$
\State \Return $\{\hat{d}, \hat{c}, \hat{r}, \rho\}$
\end{algorithmic}
\end{algorithm}

\section{Theoretical Guarantees and FLOP Scaling}

\begin{theorem}[Exact $2^\phi$ FLOP Doubling Rate]
Under the constraint $\alpha \cdot \beta^2 \cdot \gamma^2 = 2$, scaling a convolutional network by compound factor $\phi \ge 0$ scales total floating-point operations by exactly:
\begin{equation}
\text{FLOPs}(\mathcal{N}_{\text{scaled}}) = 2^\phi \cdot \text{FLOPs}(\mathcal{N}_{\text{base}})
\end{equation}
\end{theorem}

\begin{proof}
The total FLOPs of a convolutional network scales proportionally to depth $d$, quadratically with channel width $w$ (since conv layers compute $C_{\text{in}} \times C_{\text{out}}$), and quadratically with image resolution $r$ (pixel count $H \times W$):
\begin{equation}
\text{FLOPs} \propto d \cdot w^2 \cdot r^2
\end{equation}
Substituting the compound scaling rules $d = \alpha^\phi, w = \beta^\phi, r = \gamma^\phi$:
\begin{equation}
\text{FLOPs}_{\text{scaled}} = (\alpha^\phi) \cdot (\beta^\phi)^2 \cdot (\gamma^\phi)^2 \cdot \text{FLOPs}_{\text{base}} = (\alpha \beta^2 \gamma^2)^\phi \cdot \text{FLOPs}_{\text{base}}
\end{equation}
Substituting the normalization constraint $\alpha \beta^2 \gamma^2 = 2$:
\begin{equation}
\text{FLOPs}_{\text{scaled}} = 2^\phi \cdot \text{FLOPs}_{\text{base}}
\end{equation}
proving that every unit increment in $\phi$ doubles the required compute budget uniformly across all network dimensions.
\end{proof}

\subsection{Limitations}
\begin{itemize}
    \item \textbf{Memory Footprint}: Scaling resolution $r$ increases GPU activation memory by $r^2$, which can trigger out-of-memory errors during training faster than depth or width scaling alone.
\end{itemize}

\section{Complexity Analysis}

\subsection{Time Complexity}
\begin{itemize}
    \item \textbf{Stage Parameter Derivation}: Traverses $S$ network stages in $\mathcal{O}(S)$ steps.
    \item \textbf{Total Time Complexity}: $\Theta(S)$ operations.
\end{itemize}

\subsection{Space Complexity}
\begin{itemize}
    \item \textbf{Storage}: $\mathcal{O}(S)$ memory for scaled depth and channel lists.
\end{itemize}

\section{Exactness and Error Analysis}

The scaling rules execute exact floating-point power equations followed by deterministic integer hardware divisibility rounding.

\section{Worked Numerical Example}

Let $S = 2, d = [1, 2], c = [16, 32], r_0 = 224, \phi = 1.0, \alpha = 1.2, \beta = 1.1, \gamma = 1.15, \delta = 8$.

\begin{enumerate}
    \item Multipliers ($\phi = 1.0$):
    $d_{\text{mult}} = 1.2, w_{\text{mult}} = 1.1, r_{\text{mult}} = 1.15$.
    \item \textbf{Scaled Depths}:
    $\hat{d}_1 = \lceil 1 \times 1.2 \rceil = 2, \hat{d}_2 = \lceil 2 \times 1.2 \rceil = \lceil 2.4 \rceil = 3 \implies \hat{d} = [2, 3]$.
    \item \textbf{Scaled Channels} (aligned to $\delta = 8$):
    Stage 1: $16 \times 1.1 = 17.6 \to \text{round}(17.6 / 8) \times 8 = 16$. (Since $16 < 0.9(17.6) = 15.84$ is false, $\hat{c}_1 = 16$).
    Stage 2: $32 \times 1.1 = 35.2 \to \text{round}(35.2 / 8) \times 8 = 32$. (Check: $32 < 0.9(35.2) = 31.68$ is false, $\hat{c}_2 = 32$).
    \item \textbf{Scaled Resolution}:
    $\hat{r} = \text{round}(224 \times 1.15) = 258$.
    \item \textbf{FLOP Ratio}:
    $\rho = 1.2 \times (1.1)^2 \times (1.15)^2 = 1.2 \times 1.21 \times 1.3225 \approx 1.92$.
\end{enumerate}

\section{Numerical Considerations and Edge Cases}

\begin{itemize}
    \item \textbf{Hardware Alignment}: Rounding channels to multiples of 8 or 16 maximizes memory coalescing and GPU tensor core GEMM efficiency.
\end{itemize}

\begin{thebibliography}{9}
\bibitem{tan2019efficientnet} M.~Tan and Q.~V.~Le, ``EfficientNet: Rethinking Model Scaling for Convolutional Neural Networks,'' in \emph{International Conference on Machine Learning (ICML)}, pp.~6105--6114, 2019.
\bibitem{tan2021efficientnetv2} M.~Tan and Q.~V.~Le, ``EfficientNetV2: Smaller Models and Faster Training,'' in \emph{International Conference on Machine Learning (ICML)}, pp.~10096--10106, 2021.
\end{thebibliography}

\end{document}
